% Part: lambda-calculus
% Chapter: introduction
% Section: reduction

\documentclass[../../../include/open-logic-section]{subfiles}

\begin{document}

\olfileid{lam}{int}{red}
\olsection{Reduksi Suku Lambda}

Apa kang bisa ditindakake marang suku lambda?
Disederhanakake. Yen $M$ lan $N$ iku
sembarang suku lambda lan $x$ sembarang variabel,
kita bisa nggunakake $\Subst{M}{N}{x}$ kanggo
nyebut asil nyelehake $N$ minangka pangganti~$x$ ing~$M$,
sawise ngganti jeneng variabel kaiket ing $M$
kang bakal ngganggu variabel bebas ing~$N$
sawise substitusi. Umpamane,
\[
\Subst{(\lambd[w][xxw])}{yyz}{x} = \lambd[w][(yyz)(yyz)w].
\]

\begin{digress}
Notasi liya kanggo substitusi yaiku $[N/x]M$,
$[x/N]M$, lan uga $M[x/N]$. Ati-ati!
\end{digress}

Miturut intuisi, $(\lambd[x][M])N$ lan
$\Subst{M}{N}{x}$ nduweni teges kang padha;
tumindak ngganti suku kapisan nganggo suku
kapindho diarani \emph{kontraksi~$\beta$}.
$(\lambd[x][M])N$ diarani \emph{redex},
dene $\Subst{M}{N}{x}$ diarani
\emph{kontraktume}. Umume, yen suku $P$
bisa diowahi dadi~$P'$ kanthi kontraksi~$\beta$
ing sawijining subsuku, kita nyebut yen $P$
\emph{direduksi-$\beta$ dadi $P'$ sajroning
siji langkah}, lan nulis $P \redone P'$.
Yen saka $P$ kita bisa ngasilake~$P'$
kanthi sawetara reduksi siji langkah
(kalebu ora ana langkah), mula $P$
\emph{direduksi-$\beta$} dadi~$P'$;
kanthi simbol, $P \red P'$. Suku kang
ora bisa direduksi-$\beta$ maneh diarani
\emph{ora bisa direduksi-$\beta$},
utawa \emph{normal-$\beta$}. Kita bakal
nyebut ``direduksi'' tinimbang
``direduksi-$\beta$'', lan sateruse,
yen kontekse wis cetha.

Gatekna sawetara tuladha.
\begin{enumerate}
\item Kita duwe
\begin{align*}
(\lambd[x][xxy]) \lambd[z][z] & \redone (\lambd[z][z])(\lambd[z][z]) y \\
& \redone (\lambd[z][z]) y \\
& \redone y.
\end{align*}
\item ``Nyederhanakake'' suku bisa ndadekake
  suku iku luwih ruwet:
\begin{align*}
(\lambd[x][xxy])(\lambd[x][xxy]) & \redone (\lambd[x][xxy])(\lambd[x][xxy])y \\
& \redone (\lambd[x][xxy])(\lambd[x][xxy])yy \\
& \redone \dots
\end{align*}
\item Reduksi uga bisa ninggalake suku tanpa
  owah-owahan:
\[
(\lambd[x][xx])(\lambd[x][xx]) \redone (\lambd[x][xx])(\lambd[x][xx]).
\]
\item Kajaba iku, sawetara suku bisa direduksi
  kanthi luwih saka siji cara; umpamane,
\[
(\lambd[x][(\lambd[y][yx]) z)] v \redone (\lambd[y][yv]) z
\]
kanthi ngontraksi aplikasi kang paling njaba; lan
\[
(\lambd[x][(\lambd[y][yx]) z)] v \redone (\lambd[x][zx]) v
\]
kanthi ngontraksi aplikasi kang paling njero.
Nanging gatekna yen ing kasus iki, loro suku
iku banjur direduksi maneh dadi suku kang
padha, yaiku $zv$.
\end{enumerate}

Asil pungkasan ing tuladha pungkasan iku
dudu mung kabeneran, nanging nuduhake
sipat kalkulus lambda kang jero lan wigati,
kang diarani ``sipat Church--Rosser''.

\end{document}
